% Rectificación quirúrgica: generación conjunta, profundidad arbitraria y
% estatuto tipado de los transportes ternarios.  Este bloque pertenece al
% capítulo 27 y precede a la separación expositiva de los cuatro lectores.

\section{Teorema conjunto de generación coinductiva y profundidad arbitraria}
\label{sec:c27-generacion-coinductiva-profundidad-arbitraria}

La estructura de partida de este teorema no es el conjunto desnudo de nueve
símbolos. Es el sistema tipado
\[
 \mathfrak H_9=
 \bigl(\mathrm{APP},\mathrm{TRIT},\mathrm{TPK};
       \widehat{\mathcal X}^{\mathrm{enr}},\mathcal U,
       \Gamma_9,\operatorname{Ext},\operatorname{tr}\bigr),
\]
construido sobre la doble realización posicional y aritmética de
\(\mathcal D_9=\{1,\ldots,9\}\). APP acumula y pliega mediante sus hojas
\(\Sigma\) y \(\Pi\), conservando residuo y cociente; TRIT orienta cada
transición en uno de sus tres regímenes; y TPK selecciona, transporta,
memoriza y retorna sobre la fibra enriquecida. En particular,
\[
 M_{\rm ph}:\mathcal D_9\longrightarrow\{+1,-1,0\},
 \qquad
 m_{\rm ph}:=\bigl(M_{\rm ph}(1),\ldots,M_{\rm ph}(9)\bigr)^{\mathsf T},
\]
\[
 \operatorname{Tra}_t(y,d):=\operatorname{Lift}(T_d)(y),
 \qquad
 \operatorname{Upd}_t:=\operatorname{Rec}_t\circ\operatorname{Mem}_t,
\]
y, para todo \(x\in\widehat{\mathcal X}^{\mathrm{enr}}\),
\[
 \mathcal U_t(x)
 =\operatorname{Upd}_t\!\left(
   \operatorname{Tra}_t\!\left(
    \operatorname{Sel}_t\!\left(x,M_{\rm ph}(g(t))\right),d(t)
   \right)\right).
\]
Esta composición es el dato dinámico efectivo: la carta de nueve símbolos
es su soporte finito, no un sustituto de la dinámica. En particular,
\(M_{\rm ph}(g(t))\) es el valor escalar que recibe el selector, mientras
que \(m_{\rm ph}\) es el vector de los nueve valores de la función; ninguno
se compone directamente con un operador de estado.

Para \(r\in6\mathbb N_{>0}\), sean
\(\widehat{\mathcal X}^{\mathrm{enr}}_r\) los estados truncados con hoja,
orientación, residuo, cociente, acarreo, ruta, frontera y memoria, y sean
\[
 \operatorname{Ext}_{r+6,r}\subseteq
 \widehat{\mathcal X}^{\mathrm{enr}}_r\times
 \widehat{\mathcal X}^{\mathrm{enr}}_{r+6},
 \qquad
 \operatorname{tr}_{r+6,r}:
 \widehat{\mathcal X}^{\mathrm{enr}}_{r+6}
 \longrightarrow
 \widehat{\mathcal X}^{\mathrm{enr}}_r
\]
la relación de prolongación admisible y el truncamiento. Para
\(x_r\in\widehat{\mathcal X}^{\mathrm{enr}}_r\), la fibra local
\[
 \operatorname{Ext}_{r+6,r}(x_r)
 :=\{x_{r+6}:(x_r,x_{r+6})\in\operatorname{Ext}_{r+6,r}\}
\]
puede contener varios hijos. La compatibilidad correcta es
\[
 (x_r,x_{r+6})\in\operatorname{Ext}_{r+6,r}
 \Longrightarrow
 \operatorname{tr}_{r+6,r}(x_{r+6})=x_r.
\]
En \(R_{36}\), los tres renglones recuperados
\(222220,021222,102011\), transportados al calibre activo por la permutación
\(p_0=(0,1,2,4,5,3)\) en índices desde cero ---equivalentemente,
\(p_1=(1,2,3,5,6,4)\) en índices desde uno---, portan antes de toda
realización real los caracteres
internos de clausura, propagación y autoescala. Su triple, junto con el
registro firmado \(U_{12}\) y el sello \(K\), porta la coordenada interna de
\(\alpha\). Para cada
\(\chi\in\{\pi,\varphi,e,\alpha\}\), escribimos el estado así enriquecido
como
\[
 x_r^\chi=(\widetilde x_r,\chi_{\rm HMT},
            \varepsilon_r^\chi,N_r^\chi).
\]
Aquí \(\varepsilon_r^\chi\) es el residuo interno producido por el cociente
APP en el estado, no \(\operatorname{frac}(P_{\rm ar}(\chi))\). La división
euclídea interna fija
\[
 d_r^\chi=\operatorname{Dig}_{729}(\varepsilon_r^\chi),
 \quad 0\le d_r^\chi<729,
 \quad
 \varepsilon_{r+6}^\chi=729\varepsilon_r^\chi-d_r^\chi,
 \quad
 N_{r+6}^\chi=729N_r^\chi+d_r^\chi.
\]
La selección, el transporte y la actualización sobre la sección
\(\chi\) se tipan por
\[
 \operatorname{Sel}_r^\chi(x)
 :=\operatorname{Sel}\!\left(
 x,M_{\rm ph}(\rho_9(r/6+1))\right),
 \qquad
 \operatorname{Tra}_r^\chi(y,d):=\operatorname{Lift}(T_d)(y),
 \qquad
 \operatorname{Upd}_r^\chi:=\operatorname{Rec}_r^\chi\circ
 \operatorname{Mem}_r^\chi.
\]
La transición completa es
\[
 \mathcal U_r^\chi(x)
 :=\operatorname{Upd}_r^\chi\!\left(
  \operatorname{Tra}_r^\chi\!\left(
   \operatorname{Sel}_r^\chi(x),d_r^\chi
  \right)\right),
\]
y sólo lee el estado enriquecido anterior: fase, memoria, acarreo, hoja,
orientación, ruta, frontera, firma, registro, carácter y residuo internos.
No lee una cifra futura, \(P_{\rm ar}\), ni una horquilla de Machin,
factorial o Perron. Definimos la subrelación forward
\[
 \operatorname{Ext}^{\chi,\rightarrow}_{r+6,r}
 :=\{(x_r^\chi,\mathcal U_r^\chi(x_r^\chi))\}
 \subseteq\operatorname{Ext}_{r+6,r}.
\]
No es la selección retrospectiva
\(\operatorname{Ext}\cap\{P_{\rm ar}=\chi\}\): el carácter ya es una
coordenada causal de \(x_{36}^\chi\). Las ecuaciones anteriores dan
\(\lfloor N_{r+6}^\chi/729\rfloor=N_r^\chi\) y conservan hoja, orientación,
acarreo, frontera, ruta, firma, registro y genealogía.

La relación completa \(\operatorname{Ext}\) sigue siendo multivaluada. La
unicidad requerida es la de cada sección forward sobre su semilla R36
producida,
\[
 \#\Bigl\{(x_r^\chi)_{r\ge36}:x_{36}^\chi\text{ fijado},\ 
 (x_r^\chi,x_{r+6}^\chi)
 \in\operatorname{Ext}^{\chi,\rightarrow}_{r+6,r}\Bigr\}=1,
\]
no la igualdad falsa \(\#\operatorname{Ext}_{r+6,r}(x_r)=1\) en cada puerta
local. Esta compatibilidad se prolonga a todo nivel. La notación
\[
 w_6\longrightarrow w_{12}\longrightarrow w_{18}\longrightarrow w_{24}
 \longrightarrow w_{30}\longrightarrow R_{36}\longrightarrow G_9
\]
designa estas aplicaciones y sus cambios de tipo, no una lista de rótulos.
La holonomía \(\Gamma_9\) satisface
\[
 g(k+9)=g(k),
 \qquad q^{(9)}(k+9)=q^{(9)}(k)+1,
\]
de modo que retorna la fase mientras avanzan el cociente, la ruta y la
memoria del estado.

Una historia canónica genérica basada en \(a_{(1,1)}\) y
\(t_j=M_{\rm ph}(\rho_9(j+1))\) acredita la no vacuidad de
\(\operatorname{Ext}\); no individualiza por sí sola las cuatro secciones
anteriores. Éstas quedan empalmadas por sus coordenadas internas R36 y por la
regla forward que acaba de definirse.

\begin{theorem}[Construcción inaugural desde las semillas]
\label{thm:c27-construccion-inaugural-semillas}
Sea
\[
 I_9=\{0,\ldots,8\},\qquad D=\{N,E,S,O\},\qquad
 \mathcal S=I_9^2\times D,
\]
y sea \(\Omega_{\mathrm{seed}}=\mathcal S_+\times\mathcal S_\times\) el
dominio de los dos cursores que realizan las hojas aditiva y multiplicativa
de APP. Entonces
\[
 |\Omega_{\mathrm{seed}}|=(9^2\cdot4)^2=104\,976,
\]
y la recurrencia de \(54\) pasos y la reducción ternaria definen
\[
 \Omega_{\mathrm{seed}}
 \xrightarrow{\ T_{54}^{\rm seed}\ }
 \mathcal U_6^{\mathrm{cat}}
 \xrightarrow{\ \rho_3\ }
 \mathcal W_6
 \hookrightarrow\mathbb F_3^6,
\]
con
\[
 |\mathcal U_6^{\mathrm{cat}}|=468,
 \qquad |\mathcal W_6|=243,
 \qquad |\mathbb F_3^6|=729.
\]
Las fibras satisfacen exactamente
\[
 432\cdot144+18\cdot432+18\cdot1944=104\,976.
\]
La relación de extensión del TPK prolonga las palabras supervivientes a la
familia de secciones compatibles de profundidad arbitraria. Las subrelaciones
forward ancladas en R36 individualizan la generación coinductiva
correlacionada de \(\pi,\varphi,e\) y \(\alpha\); la relación completa puede
seguir teniendo varias prolongaciones locales. Sobre esa familia se
construyen, conjuntamente, la sombra arquimediana y la realización
excepcional del \cref{thm:c27-tesis-rectora-doble-proyeccion}.
\end{theorem}

\begin{proof}
El producto cartesiano anterior enumera exhaustivamente las posiciones y
orientaciones de ambos cursores. La aplicación \(T_{54}^{\rm seed}\) se obtiene al
componer en cada paso la selección de hoja, el transporte cardinal, la
actualización del cociente y la memoria de acarreo. La enumeración de su
imagen contiene \(468\) sextetos distintos; la suma de las cardinalidades de
sus fibras es la identidad expuesta. La aplicación componente a componente
\(\rho_3\) contiene \(243\) palabras distintas. Finalmente, la no vacuidad en
todo horizonte y la compatibilidad de \(\operatorname{Ext}\) con los
truncamientos producen el límite inverso. La funcionalidad de cada
\(\operatorname{Ext}^{\chi,\rightarrow}\), a partir del estado R36 que ya
porta \(\chi_{\rm HMT}\), demuestra por inducción la unicidad de esa sección
global sin exigir que cada fibra local de \(\operatorname{Ext}\) sea
unitaria. Cada afirmación es una identidad interna o una consecuencia de la
compatibilidad proyectiva; ninguna utiliza como entrada un número real, una
constante reconocida o una tabla decimal.
\end{proof}

\begin{theorem}[Tesis rectora: un continuo y dos proyecciones]
\label{thm:c27-tesis-rectora-doble-proyeccion}
Sea \(\Omega_\Gamma^{\mathrm{cal}}\subset
X_\infty^{\mathrm{cal}}\) el espacio de secciones supervivientes de las
vueltas enriquecidas del TPK. El acarreo de las tres vueltas
\(\gamma_-,\gamma_0,\gamma_+\) define, sin evaluación arquimediana previa,
el homeomorfismo
\[
 \beta:\Omega_\Gamma^{\mathrm{cal}}\xrightarrow{\ \sim\ }
 (\mathbb F_3^6)^{\mathbb N}=: \mathcal W_\infty.
\]
Las formas equilibradas normalizadas
\[
 a_k+c_k=\rho_k+3c_{k+1},\qquad
 \rho_k\in\{-1,0,+1\},
\]
construyen el anillo ordenado \(\mathbb Z_{\rm HMT}\), su cuerpo de
fracciones \(\mathbb Q_{\rm HMT}\) y la completación
\[
 \mathbb R_{\rm HMT}:=
 \operatorname{Cau}(\mathbb Q_{\rm HMT})/{\asymp}.
\]
La estructura discreta conjunta es el producto fibrado
\[
\mathcal C_{\mathrm{cont}}^{\mathrm{disc}}
=\bigl(\mathbb Z_{\rm HMT}\times
\Omega_\Gamma^{\mathrm{cal}}\bigr)
\times_{\mathcal W_\infty}X_\infty^{\mathrm{cal}}.
\]
Sus dos aplicaciones naturales son
\[
 P_{\mathrm{ar}}:\mathbb Z_{\rm HMT}\times
 \Omega_\Gamma^{\mathrm{cal}}\longrightarrow\mathbb R_{\rm HMT},
 \qquad
 Q_\infty:X_\infty^{\mathrm{cal}}\longrightarrow
 (\mathscr C_W^{\mathrm{cal}})^{\mathbb N}.
\]
Si \(\beta(\xi)=(b_q)_{q\geq0}\),
\(\nu(b)=\sum_{i=1}^{6}b_i3^{6-i}\) y
\[
 q_n(m,\xi)=m+\sum_{q=0}^{n-1}
 \frac{\nu(b_q)}{729^{q+1}},
\]
entonces \((q_n)\) es de Cauchy y \(P_{\mathrm{ar}}(m,\xi)\) es su clase
en \(\mathbb R_{\rm HMT}\). Los cilindros racionales
\[
 I_n(m;b_0,\ldots,b_{n-1})=
 \left[q_n,q_n+729^{-n}\right]
\]
son compatibles, tienen diámetro tendente a cero y la partición recurrente
en \(729\) subcilindros construye, para toda clase de Cauchy, una cinta
\((b_q)\) que la representa. En consecuencia,
\[
 P_{\mathrm{ar}}\ \text{es sobreyectiva},\qquad
 (\mathbb Z_{\rm HMT}\times\Omega_\Gamma^{\mathrm{cal}})
 /{\sim_{\rm ar}}\ \cong\ \mathbb R_{\rm HMT}.
\]
Sólo después de esta construcción interna, la unicidad del cuerpo ordenado
completo arquimediano reconoce \(\mathbb R_{\rm HMT}\cong\mathbb R\).

La segunda proyección conserva la incidencia calibrada de cada nivel; en las
coordenadas \(x=((w_j,f_j))_{j\geq1}\) actúa por
\[
 Q_\infty(x)
 =\bigl(f_j,\Gamma_{f_jA_Wf_j^{-1}}(w_j)\bigr)_{j\geq1}.
\]
Así, \(\mathbb R_{\rm HMT}\), reconocido posteriormente como \(\mathbb R\),
es la sombra arquimediana obtenida al contraer la genealogía completa de
cilindros. La proyección excepcional conserva esa genealogía como incidencia
de frontera y admite la realización sucesiva
Paley--Witt--Golay--Leech--\(V^\natural\)--Monstruo. Las dos ramas proceden
del mismo estado y ninguna se reconstruye retrospectivamente desde la otra.
\end{theorem}

\begin{proof}
Las tres vueltas enriquecidas existen después de cada prefijo y su acarreo
las distingue por \(-1,0,+1\). La concatenación y el truncamiento son
inversos, lo que prueba la biyectividad y continuidad de \(\beta\). Para una
clase de Cauchy de \(\mathbb Q_{\rm HMT}\), se elige primero su cilindro
unitario ordenado y después, recursivamente, el único subcilindro semiabierto
de la partición interna en \(729\) piezas que contiene la clase; en una
frontera común, las dos escrituras se identifican por \(\sim_{\rm ar}\).
La sucesión de extremos izquierdos es precisamente \((q_n)\), converge a la
clase dada y prueba la sobreyectividad sin presuponer una expansión de
\([0,1]\) ni un punto de \(\mathbb R\). La identificación con el cuerpo real
es el teorema de unicidad aplicado después a \(\mathbb R_{\rm HMT}\).

Los subespacios que satisfacen el cierre de un carácter concreto son filtros
de supervivencia dentro de la capacidad global. Su función es individualizar
las secciones de \(\pi,e,\varphi,\alpha\); la sobreyectividad del continuo
procede de la completación de \(\mathbb Q_{\rm HMT}\). Esta separación de
dominios impide
convertir una condición de selección de constantes en una premisa de la
cobertura de \(\mathbb R\).

En la segunda rama, la compatibilidad
\(r_{n,m}Q_n=Q_m p_{n,m}\) permite pasar al límite y define \(Q_\infty\).
La conjugación por \(f_j\) transporta el calibre sin borrar la incidencia.
Los funtores posteriores de código, retículo y álgebra de vértices realizan
la cadena excepcional. La igualdad de las cintas en el producto fibrado
construye la fuente conjunta y conserva informaciones complementarias: valor
arquimediano en la primera; incidencia, frontera y genealogía en la segunda.
\end{proof}

\begin{corollary}[Expansión infinita como objeto matemático]
\label{cor:c27-expansion-infinita-objeto}
Para cada carácter numérico \(\chi\) publicado por el estado coinductivo, la
lectura decimal es una secuencia completa
\[
 \operatorname{Dec}_\chi(x_\infty)
 =(d_{\chi,1},d_{\chi,2},\ldots)
 \in\{0,\ldots,9\}^{\mathbb N}.
\]
Su prefijo de longitud \(N\) es la proyección
\(\rho_N\operatorname{Dec}_\chi(x_\infty)\). Por ello, mil, diez mil o un
millón de cifras son cortes del mismo objeto infinito; la generación no
consiste en elevar sucesivamente una cota experimental.
\end{corollary}

\begin{theorem}[Generación coinductiva correlacionada de \(\pi\), \(\varphi\), \(e\) y \(\alpha\)]
\label{thm:c27-cuadrirrelacion-generada}
El perfil completo del sistema \(\mathfrak H_9\) determina la familia
correlacionada de secciones forward compatibles
\[
 \boldsymbol x_\infty=
 (x_\infty^\pi,x_\infty^\varphi,x_\infty^e,x_\infty^\alpha)
 \in
 \widehat\Omega_\pi\times_{R_{36}}
 \widehat\Omega_\varphi\times_{R_{36}}
 \widehat\Omega_e\times_{R_{36}}\widehat\Omega_\alpha,
 \qquad
 \widehat\Omega_\chi:=
 \varprojlim_{r\ge36}
 (\widehat{\mathcal X}^{\mathrm{enr},\chi}_r,
  \operatorname{Ext}^{\chi,\rightarrow}_{r+6,r}),
\]
y cuatro caracteres correlacionados
\[
 \bigl(\Pi_{\mathrm{cl}},\Pi_{\mathrm{aut}},
       \Pi_{\mathrm{prop}},\Pi_{\mathrm{dod}}\bigr)
       (\boldsymbol x_\infty)
 =\bigl(\pi_{\mathrm{HMT}},\varphi_{\mathrm{HMT}},
        e_{\mathrm{HMT}},\alpha_{\mathrm{HMT}}\bigr).
\]
Los tres primeros caracteres proceden, respectivamente, de la única
órbita de clausura, la única órbita de autoescala y la única
órbita de propagación del catálogo enriquecido. El cuarto procede de la
coordenada primitiva firmada de la familia terminal correlacionada y de su
cierre dodecafásico. Las cuatro coordenadas son, por tanto, salidas de la
generación coinductiva correlacionada de profundidad arbitraria gobernada por
la misma dinámica APP--TRIT--TPK, con retorno de fase, avance de memoria y
conservación genealógica, aunque sus aplicaciones de publicación sean
distintas.
\end{theorem}

\begin{proof}
El censo exhaustivo
\(104,976\to468\to243\subset\mathbb F_3^6\) descompone la imagen realizada
en \(43\) órbitas diedrales. Los invariantes de fibra dejan una sola órbita
en cada una de las clases de clausura, propagación y autoescala; el calibre
interno orienta en ellas, sin consulta decimal, los representantes
\[
 w_6^{\mathrm{cl}}=010211,
 \qquad w_6^{\mathrm{prop}}=201101,
 \qquad w_6^{\mathrm{aut}}=121200.
\]
Las subrelaciones forward conservan el carácter interno y los predicados que
definen las tres clases; Machin, la recurrencia factorial y Perron no eligen
ningún hijo de \(\operatorname{Ext}\). En el
canal de clausura, la composición pentarregional con retorno unitario produce
las razones \(1/5\), \(1/239\) y la relación \(120/119\); la identidad de
Machin reconoce exactamente la media vuelta. En el canal de propagación, la
ley semigrupal normalizada fuerza
\((n+1)a_{n+1}=a_n\), luego \(a_n=1/n!\) y \(P_1=e\). En el canal de
autoescala, la incidencia
\[
 F_{\mathrm{av}}=\begin{pmatrix}0&1\\1&1\end{pmatrix}
\]
posee un único rayo positivo y su ecuación propia fuerza
\(x^2-x-1=0\), \(x>0\), esto es, \(x=\varphi\).

En el perfil TPK completo, la lectura terminal
\(x_{\mathrm{term}}\mapsto(B_{\mathrm{term}},U_{12}^{\mathrm{sgn}})\)
publica simultáneamente la carta de los tres canales y la coordenada
primitiva firmada. La
segunda se identifica reversiblemente con el sello \(K\) y sus diferencias
dodecafásicas. La recurrencia euclídea con acarreo
\[
 s_m=P_m+E_m-\Phi_m-K_m+c_{m+1},\qquad
 A_m=s_m\bmod1000,\qquad
 c_m=(s_m-A_m)/1000
\]
produce \(\alpha_{\mathrm{HMT}}\) sin recibirla como entrada. La
caracterización analítica independiente comienza con el núcleo finito
\(E_{21/22}^{(9)}\), que determina una raíz simple \(\xi_9\) en
\((0,10^{-2})\) sin consultar la salida dodecafásica. La iteración del mismo
lector graduado produce el sistema inverso
\(E_{21/22}^{(6)}\leftarrow E_{21/22}^{(9)}\leftarrow
E_{21/22}^{(12)}\leftarrow\cdots\). Sus truncamientos conmutan con los de la
vía dodecafásica y ambas familias determinan en cada profundidad el mismo
cilindro superviviente. Como sus diámetros tienden a cero, la intersección es
unitaria y se obtiene exactamente
\(\alpha_A=\alpha_B=\alpha_{\mathrm{HMT}}\). La igualdad
\(\operatorname{Tr}_9(\xi_9^{-1})=
\operatorname{Tr}_9(\alpha_{12}^{-1})\) es únicamente el certificado finito
de orden nueve de esa identidad de secciones completas. Así, la separación
de lectores conserva el origen común de las cuatro coordenadas sin introducir
una prolongación pendiente ni un valor objetivo.
\end{proof}

\begin{proposition}[Registro dodecafásico del estado completo anterior a \(\alpha\)]
\label{prop:c27-registro-u-k-producido}
El estado terminal enriquecido posee la doble publicación
\[
x_{\mathrm{term}}
\xrightarrow{(\pi_{\mathrm{term}},R_{\mathrm{sgn}})}
(B_{\mathrm{term}},U),
\]
donde
\[
U=(2378,1406,2479,-452,998,-551,
-668,-204,-371,-322,-28,-997).
\]
Al ordenar las coordenadas por las tres órbitas de paso tres, la transformada
\[
\mathcal H_{12}
=P_{\mathrm{orb}}^{-1}(H_4\oplus H_4\oplus H_4)P_{\mathrm{orb}},
\qquad H_4^2=4I_4,
\]
posee inversa \(\mathcal H_{12}/4\) sobre su imagen integral y determina
unívocamente
\[
K=(234,543,140,729,659,824,621,058,914,794,146,601).
\]
Por tanto, dentro del perfil TPK completo,
\[
x_{\mathrm{term}}\xrightarrow{R_{\mathrm{sgn}}}U
\xrightarrow{\mathcal H_{12}^{-1}}K
\]
es una composición exacta anterior al acarreo que publica \(\alpha\). La matriz
visible \(B_{\mathrm{term}}\) y el registro firmado \(U\) son dos
coordenadas del mismo estado; la segunda conserva los campos de orientación,
oposición, cociente y registro que la primera contrae. La composición
\[
 \Omega_{\mathrm{seed}}\longrightarrow
 \mathcal U_6^{\mathrm{cat}}\longrightarrow
 \widehat{\mathcal X}^{\mathrm{enr}}_{\mathrm{term}}
 \xrightarrow{R_{\mathrm{sgn}}}U
 \xrightarrow{\mathcal H_{12}^{-1}}K
\]
construye cada coordenada por APP--TRIT--TPK y conserva sus registros. Sus
entradas son las semillas, las hojas y las reglas operatorias; \(U\), \(K\),
sus diferencias, \(\alpha\) y las expansiones publicadas pertenecen al
codominio de etapas sucesivas.
\end{proposition}

\begin{theorem}[Clausura bidireccional de la segunda vía de \(\alpha\)]
\label{thm:c27-alpha-segunda-via-bidireccional}
La prolongación del núcleo de vacancias hasta los grados \(7{:}9\) determina
\[
\begin{aligned}
B_9(\alpha):={}&\log_{10}\!\left(\frac{10}{9}\right)
+\frac74\alpha^2-\frac{\alpha^4}{20}
+\frac{2\alpha^5}{21}+\frac{\alpha^6}{46}\\
&+\frac{\alpha^7}{120}-\frac{\alpha^8}{45}
-\frac{2\alpha^9}{495}.
\end{aligned}
\]
Al elevar el giro a la incógnita positiva \(T\), el resolvente equivale a
\[
\boxed{\alpha T^2-B_9(\alpha)T+\alpha^3=0.}
\]
Con
\(\xi=\operatorname{arcosh}(B_9(\alpha)/(2\alpha^2))\), las soluciones son
\[
T_\pm=\alpha e^{\pm\xi},
\qquad T_+T_-=\alpha^2,
\qquad
J_\alpha(T)=\frac{\alpha^2}{T},
\qquad J_\alpha(T_\pm)=T_\mp.
\]
La cámara \(6<T<7\) selecciona \(T_+\); la clausura finita
\(\pi_9=T_+/2\) se prolonga mediante \(\Gamma_9\) hasta
\(\pi_{\mathrm{HMT}}\). La identidad
\(\alpha=\sqrt{T_+T_-}\) fija la media geométrica de las ramas conjugadas.
El jet \(7{:}9\), la raíz simple, la involución y el retorno constituyen la
segunda vía interna completa.
\end{theorem}

\begin{theorem}[Acoplamiento exacto del sector \(\alpha\) con el Monstruo]
\label{thm:c27-alpha-rama-excepcional}
En la carta incidencial del mismo registro dodecafásico, las raíces radiales
de \(A_2^{12}\) determinan
\[
v_\alpha=4\rho_1+\rho_2+\cdots+\rho_{12},
\qquad v_\alpha^2=54,
\]
y el vecino de índice tres
\[
N_{\alpha,0}
=\{x\in N(C_W):\langle x,v_\alpha\rangle\equiv0\pmod3\},
\qquad
N_\alpha=N_{\alpha,0}+\mathbb Z\frac{v_\alpha}{3}.
\]
El retículo \(N_\alpha\) es par, unimodular, de rango veinticuatro y sin
raíces, de modo que \(N_\alpha\cong\Lambda_{24}\). La cadena
\[
C_W\longrightarrow N(A_2^{12})
\xrightarrow{v_\alpha}\Lambda_{24}
\longrightarrow V^\natural
\xrightarrow{\operatorname{Aut}}\mathbb M
\]
expone la bisagra excepcional. El escalar \(\alpha_{\mathrm{HMT}}\) y el
vector \(v_\alpha\) pertenecen a codominios distintos y comparten el estado
dodecafásico de origen; esta correlación tipada es el enganche con las
simetrías del Monstruo.
\end{theorem}

\begin{proposition}[Estatuto de los transportes \(L_0,L_1\)]
\label{prop:c27-estatuto-l0-l1}
Los endomorfismos \(L_0,L_1\in\operatorname{End}(\mathbb F_3^6)\) y el
calendario \((L_0,L_0,L_1,L_1)\) pertenecen a la ley de transporte del estado
TPK enriquecido. Las seis transiciones publicadas de rango completo fijan
cada matriz de manera única por \(L=X^{-1}Y\) en \(\mathbb F_3\), y el
calendario es el único de los \(16\) calendarios binarios que conserva
simultáneamente las tres biografías estructurales. Su función es prolongar
los representantes ya seleccionados; la evaluación arquimediana y la
nomenclatura convencional se aplican después.
\end{proposition}

El dominio de esta proposición es \(\mathfrak H_9\), el sistema completo
APP--TRIT--TPK con estado enriquecido. El censo APP constituye su etapa
inicial; TRIT aporta la orientación y TPK construye la prolongación, la
memoria y el retorno. Esta secuencia fija con exactitud el dominio de cada
flecha y la residencia de cada salida.

\begin{theorem}[Profundidad arbitraria y determinación de las expansiones completas]
\label{thm:c27-profundidad-arbitraria}
Para cada
\(\chi\in\{\pi_{\mathrm{HMT}},e_{\mathrm{HMT}},
\varphi_{\mathrm{HMT}}\}\) y cada \(N\ge1\), existe un único prefijo
decimal \(p_{\chi,N}\) de longitud \(N\), producido por un truncamiento
finito del estado enriquecido, tal que
\[
 \rho_{N+1,N}(p_{\chi,N+1})=p_{\chi,N}.
\]
La familia compatible es la familia de proyecciones de un único elemento
\[
 p_{\chi,\infty}
 \in\varprojlim_N\{0,\ldots,9\}^{N}
 \cong\{0,\ldots,9\}^{\mathbb N},
\]
y la intersección de sus cilindros arquimedianos consta de un único real.
Por consiguiente, la conexión genera la expansión completa; cada precisión
finita es una proyección de esa salida infinita.
\end{theorem}

\begin{proof}
La prueba tiene dos etapas que no pueden invertirse. En la primera, enteramente
anterior a la carta decimal, la sección superviviente de
\(\mathfrak H_9\) publica tres caracteres exactos. El carácter de clausura
queda fijado por la pentafibra orientada, su pendiente primitiva \(1/5\), la
composición de sus cuatro compañeras \(120/119\) y el compensador entero
\(239\); el producto gaussiano correspondiente termina sobre la semirrecta
real negativa y fija unívocamente la media vuelta orientada. El carácter de
propagación es la única serie formal normalizada por
\(a_0=1\) y \((n+1)a_{n+1}=a_n\). El carácter de autoescala es el único rayo
expansivo positivo de
\[
 F_{\rm av}=\begin{pmatrix}0&1\\1&1\end{pmatrix},
 \qquad X^2-X-1=0.
\]
Estas tres condiciones son salidas del estado enriquecido; no contienen
prefijos decimales ni consultan valores convencionales.

Sólo en la segunda etapa actúa la publicación arquimediana. Para cada carácter
ya generado y cada \(N\), sea \(\operatorname{Dec}_N\) el único cilindro
decimal canónico de diámetro \(10^{-N}\) que contiene su realización exacta.
Este cilindro es una coordenada publicada del carácter: no es una fibra local
de \(\operatorname{Ext}\) ni se usa para resolverla.
La unicidad de la realización y la convención no terminal implican
\[
 \rho_{N+1,N}\operatorname{Dec}_{N+1}(\chi)
 =\operatorname{Dec}_N(\chi).
\]
Las identidades de Machin, las sumas parciales de la serie exponencial y los
cocientes consecutivos de Fibonacci--Perron proporcionan intervalos racionales
certificados para calcular \(\operatorname{Dec}_N\); no intervienen en
\(\operatorname{Ext}\), no eligen una órbita o un carácter y no seleccionan
hijos de sus fibras locales. Sí calculan el bloque coordenado de publicación,
que permanece tipado como salida posterior. Los
cilindros publicados están encajados y sus diámetros tienden a cero, de modo
que su límite inverso determina un único punto. Como \(N\) es arbitrario y
la relación estructural de extensión carece de profundidad terminal, la
conclusión vale simultáneamente para toda longitud finita y para la expansión
infinita.
\end{proof}

El portador ejecutable impone materialmente esa separación mediante tres
funciones disjuntas: primero genera los caracteres exactos, después los evalúa
y sólo al final publica bloques arquimedianos. Su propia auditoría sintáctica
falla si un evaluador entra en el generador. La ejecución de \(396\) niveles,
\(2376\) trits, \(43\) vueltas, \(387\) pares \((K,K+9)\) por canal y
\(1000\) cifras verifica una instancia material extensa del teorema; la de
\(10,000\) cifras verifica otra. Ninguna de ambas cantidades limita la
profundidad. El cuantificador \(\forall N\) y la compatibilidad
\(\rho_{N+1,N}p_{N+1}=p_N\), junto con la existencia del elemento
\(p_\infty\) del límite inverso, demuestran la salida decimal infinita; los
cálculos finitos son certificados reproducibles de sus proyecciones.

La coordenada \(\alpha_{\mathrm{HMT}}\) pertenece a la misma familia límite
correlacionada. Su vía entera se prolonga mediante
\(\mathsf T_{\alpha,\infty}\). El lector
analítico interno construye exactamente \(E_{21/22}^{(9)}\) y su raíz simple
\(\xi_9\) en el cociente de jets de orden nueve. La estabilidad comprobada a
\(10,000\) cifras certifica la propagación remota del acarreo de la vía
entera, pero es sólo una proyección finita. Los cuadrados de truncamiento de
la familia dodecafásica y de la familia de jets conmutan exactamente; la
unicidad de sus secciones compatibles identifica
\(\alpha_A=\alpha_B=\alpha_{\mathrm{HMT}}\) en
\(\mathbb R_{\rm HMT}\) antes de toda comparación metrológica. Las corridas a
nueve, mil o diez mil cifras son testigos finitos de esta identidad y no su
alcance.

\begin{proposition}[Prolongación completa de la vía entera de \(\alpha\)]
\label{prop:c27-alpha-prolongacion-completa}
Sea \(B=1000\), sean \((P_k),(E_k),(\Phi_k)\) las tres sucesiones de bloques
publicadas por el estado coinductivo y sea \((K_k)\) la sucesión
dodecafásica prolongada. Las cuatro series
\[
 p=\sum_{k\geq1}P_kB^{-k},\qquad
 e_0=\sum_{k\geq1}E_kB^{-k},\qquad
 \phi_0=\sum_{k\geq1}\Phi_kB^{-k},\qquad
 \kappa=\sum_{k\geq1}K_kB^{-k}
\]
convergen absolutamente y determinan
\[
 \alpha_A=p+e_0-\phi_0-\kappa.
\]
La recurrencia de acarreo es la escritura en base \(B\) de esta igualdad.
Para cada \(M\), un truncamiento suficientemente profundo determina las
primeras \(M\) cifras de la expansión canónica de \(\alpha_A\); el error de
la cola está acotado por una constante multiplicada por \(B^{-N}\). Por
consiguiente, \(\mathsf T_{\alpha,\infty}\) publica una sucesión decimal
completa, y sus ejecuciones a \(3000\) o \(10\,000\) cifras son proyecciones
finitas de esa salida.
\end{proposition}

\begin{proof}
Cada bloque pertenece a un conjunto entero acotado, de modo que las cuatro
series están dominadas por una serie geométrica. La suma parcial hasta \(N\)
difiere de la suma completa en \(O(B^{-N})\). La división euclídea de
derecha a izquierda transporta exactamente esa cola mediante el acarreo;
por unicidad de la expansión canónica, todo prefijo que queda separado de
una frontera de cilindro se estabiliza al aumentar \(N\), y en una frontera
se aplica la convención canónica no terminal. Esto define simultáneamente
todas las cifras, no una sucesión de objetivos numéricos independientes.
\end{proof}

\begin{corollary}[Separación entre generación y reconocimiento]
\label{cor:c27-generacion-reconocimiento}
Las identidades de Machin, la serie exponencial, el rayo de Perron y la carta
decimal son realizaciones exactas de caracteres previamente seleccionados
por \(\mathfrak H_9\). Los nombres convencionales \(\pi,e,\varphi,\alpha\)
y las tablas metrológicas reconocen las salidas; no determinan órbitas,
caracteres, transportes ni acarreos. Los prefijos que calculan son únicamente
proyecciones coordenadas posteriores de los caracteres ya generados.
\end{corollary}

\paragraph{Evaluación de los caracteres construidos.}
En la etapa de evaluación arquimediana, el carácter de propagación
determina la condición inicial \(a_0=1\) y la recurrencia
\((n+1)a_{n+1}=a_n\). El cálculo de los términos aplica esa recurrencia.
Para la autoescala, el vector fila evoluciona mediante la incidencia
previamente construida,
\[
 (x_{n+1},y_{n+1})=(x_n,y_n)F_{\rm av},
 \qquad
 F_{\rm av}=\begin{pmatrix}0&1\\1&1\end{pmatrix},
 \qquad
 \det(tI-F_{\rm av})=t^2-t-1.
\]
Los cocientes de componentes proporcionan las cotas racionales del
lector de autoescala. La implementación verifica la concordancia entre
incidencia, polinomio característico, orientación y normalización antes
de efectuar esta evaluación. Una alteración o ausencia de esos datos
interrumpe el cálculo en lugar de conservar silenciosamente el valor
anterior. Este control enlaza los caracteres ya producidos con sus
evaluadores; su lugar en la construcción es posterior al generador.
